\begin{derivation}[消去最近外部通道产生什么修正]
取 $\Omega$ 为实数以简化记号，能量相位已由基矢重定义吸收。对角能量分别为 $0,0,2\hbar\omega_{\rm r},2\hbar\omega_{\rm r}$ 的 $|0\rangle,|1\rangle,|-1\rangle,|2\rangle$，二阶 Schur 自能在中心能量处给
\begin{align*}
 (\delta H_P)_{00}&=\frac{|\hbar\Omega/2|^2}{-2\hbar\omega_{\rm r}}
 =-\frac{\hbar|\Omega|^2}{8\omega_{\rm r}},\\
 (\delta H_P)_{11}&=-\frac{\hbar|\Omega|^2}{8\omega_{\rm r}},\qquad
 (\delta H_P)_{01}=0.
\end{align*}
最近邻链中没有经过一个外部态连接 $0$ 与 $1$ 的两步路径，所以二阶非对角修正为零。两个相同的能移只给公共相位，并不在这一对称模型中改变共振条件；更高阶仍会改变振荡频率。

若外部态初始未占据，并仅保留通向最近外部态的首阶跃迁，忽略更远的 $|-2\rangle$ 对它的反馈，则
\begin{equation*}
 d_{-1}(t)\simeq-\frac{\ii\Omega^*}{2}\int_0^t
 \ee^{-2\ii\omega_{\rm r}(t-t')}d_0(t')\dd t'.
\end{equation*}
慢变 $d_0$ 的分部积分主项为
$-\Omega^*[d_0(t)-\ee^{-2\ii\omega_{\rm r}t}d_0(0)]/(4\omega_{\rm r})$，余项含 $\dot d_0$。因此外部振幅是 $O(|\Omega|/\omega_{\rm r})$，概率是其平方；方脉冲还带有开关产生的快速瞬态。这个积分同时解释了为什么仅把外部能级的能移加进两态模型，仍不能代替对实际泄漏的计算。
\end{derivation}

\begin{derivation}[虚通道链中的有效耦合与能级移动]
取正文已经定义的最近邻虚链。自由部分在格点基底中对角：
\begin{equation*}
 H_0|j\rangle=E_j|j\rangle,
\end{equation*}
相互作用只连接最近邻：
\begin{equation*}
 V
 =\frac{\hbar}{2}\sum_j
 \left(\Omega_j|j+1\rangle\langle j|
 +\Omega_j^*|j\rangle\langle j+1|\right).
\end{equation*}
目标子空间与补空间为
\begin{align*}
 P&=|0\rangle\langle0|+|N\rangle\langle N|,\\
 Q&=\sum_{j=1}^{N-1}|j\rangle\langle j|+Q_{\rm ext},
\end{align*}
其中 $Q_{\rm ext}$ 表示目标区间之外的其他格点；计算最短 $0\to N$ 路径时它们不参与最低阶非对角矩阵元。取展开能量 $E_*$，定义
\begin{equation*}
 \Delta_j=\frac{E_*-E_j}{\hbar},
\end{equation*}
则未扰动 $Q$ 预解式在内部格点上严格为
\begin{equation*}
 G_Q^{(0)}(E_*)|j\rangle
 =\frac1{E_*-E_j}|j\rangle
 =\frac1{\hbar\Delta_j}|j\rangle.
\end{equation*}

先求非对角元 $\langle N|H_{\rm eff}|0\rangle$。由于 $PVP$ 对 $N>1$ 为零，最低非零项必须含 $N$ 个最近邻顶角和 $N-1$ 个中间传播子。正文\eqnrefs{eq:virtual-path-general-series-dp85} 中对应的项是
\begin{align*}
 \Sigma_{N0}^{(N)}(E_*)
 ={}&\langle N|VQG_Q^{(0)}QVQG_Q^{(0)}\cdots
 QVQG_Q^{(0)}QV|0\rangle,
\end{align*}
其中共有 $N-2$ 个内部 $QVQ$，加上两端的 $PVQ$ 与 $QVP$，总顶角数为 $N$。在最近邻假设下，每插入一次恒等分解
\begin{equation*}
 Q=\sum_{j=1}^{N-1}|j\rangle\langle j|+Q_{\rm ext}
\end{equation*}
以后，所有不满足相邻格点选择律的项都为零，因此最低阶只剩唯一单调路径
\begin{equation*}
 0\to1\to2\to\cdots\to N-1\to N.
\end{equation*}
逐个写出矩阵元：
\begin{align*}
 \Sigma_{N0}^{(N)}(E_*)
 ={}&\langle N|V|N-1\rangle
 \frac1{E_*-E_{N-1}}
 \langle N-1|V|N-2\rangle\\
 &\times\frac1{E_*-E_{N-2}}
 \cdots
 \frac1{E_*-E_1}
 \langle1|V|0\rangle.
\end{align*}
利用
\begin{equation*}
 \langle j+1|V|j\rangle=\frac{\hbar\Omega_j}{2},
 \qquad
 E_*-E_j=\hbar\Delta_j,
\end{equation*}
得到
\begin{align*}
 \Sigma_{N0}^{(N)}(E_*)
 &=\frac{(\hbar\Omega_{N-1}/2)\cdots(\hbar\Omega_0/2)}
 {(\hbar\Delta_{N-1})\cdots(\hbar\Delta_1)}\\
 &=\frac{\hbar}{2}
 \frac{\Omega_{N-1}\Omega_{N-2}\cdots\Omega_0}
 {2^{N-1}\Delta_{N-1}\Delta_{N-2}\cdots\Delta_1}.
\end{align*}
按正文采用的二态非对角元约定
\begin{equation*}
 \langle N|H_{\rm eff}|0\rangle
 \equiv\frac{\hbar\Omega_{\rm eff}^{(N)}}2,
\end{equation*}
立即得到正文\eqnrefs{eq:nth-bragg-effective-rabi-dp83} 的首项。

为什么下一项比首项多两个顶角也可以从路径结构直接看出。最近邻格上，从 $0$ 到 $N$ 的路径净向右位移必须为 $N$。若额外增加一步向左，则为了最终仍到达 $N$，必须再增加一步向右。因此相对于最短 $N$ 顶角路径，任何更长路径至少含 $N+2$ 个顶角。若所有耦合满足 $|\Omega_j|\le\Omega_*$，所有内部能量分母满足 $|\Delta_j|\ge\Delta_*$，则下一阶非对角矩阵元的尺度至多为
\begin{equation*}
 \frac{|\delta\Sigma_{N0}|}{\hbar}
 =O\!\left(
 \frac{\Omega_*^{N+2}}{\Delta_*^{N+1}}
 \right),
\end{equation*}
对应正文\eqnrefs{eq:nth-bragg-effective-rabi-dp83} 的修正阶数；如果存在更远程连边，这个“多两个顶角”的计数不再成立，应改用一般算符范数余项界~\eqnrefs{eq:virtual-path-remainder-bound-dp85}。

再求端点的最低阶对角自能。对任一保留态 $|a\rangle\in P$，Feshbach 自能的二阶项为
\begin{equation*}
 \delta E_a^{(2)}
 =\langle a|VQG_Q^{(0)}Q V|a\rangle.
\end{equation*}
在 $Q$ 子空间插入完整本征态分解，不能预先丢掉链外邻态：
\begin{align*}
 \delta E_a^{(2)}
 &=\sum_{m\in Q}
 \langle a|V|m\rangle
 \frac1{E_*-E_m}
 \langle m|V|a\rangle\\
 &=\sum_{m\in Q}
 \frac{|\langle m|V|a\rangle|^2}{E_*-E_m}.
\end{align*}
只对非零矩阵元保留求和项，并用正文记号
\begin{equation*}
 \langle m|V|a\rangle=\frac{\hbar\Omega_{ma}}2,
 \qquad
 \Delta_m^{(a)}=\frac{E_*-E_m}{\hbar},
\end{equation*}
便有
\begin{align*}
 \delta E_a^{(2)}
 &=\sum_{m\in Q:\,\langle m|V|a\rangle\ne0}
 \frac{\hbar^2|\Omega_{ma}|^2/4}{\hbar\Delta_m^{(a)}}\\
 &=\hbar
 \sum_{m\in Q:\,\langle m|V|a\rangle\ne0}
 \frac{|\Omega_{ma}|^2}{4\Delta_m^{(a)}}.
\end{align*}
除以 $\hbar$ 即得到正文\eqnrefs{eq:nth-bragg-ac-stark-dp83}。若链在两端真正截断，则 $a=0$ 的求和只剩 $m=1$，$a=N$ 的求和只剩 $m=N-1$，从而分别约化为
\begin{equation*}
 \delta_0^{\rm AC}=\frac{|\Omega_0|^2}{4\Delta_1},
 \qquad
 \delta_N^{\rm AC}=\frac{|\Omega_{N-1}|^2}{4\Delta_{N-1}}.
\end{equation*}
若目标态位于无限最近邻格中，则还存在链外邻态，必须把相应 $m$ 项保留在同一求和式中。

最后估计中间态泄漏。对最邻近中间态，把除 $0\leftrightarrow1$ 直接耦合以外的所有项记作 $r_1$，定态方程严格写成
\begin{equation*}
 (E_*-E_1)c_1
 =\frac{\hbar\Omega_0}{2}c_0+r_1.
\end{equation*}
由 $\Delta_1=(E_*-E_1)/\hbar$ 得到恒等式
\begin{equation*}
 c_1
 =\frac{\Omega_0}{2\Delta_1}c_0
 +\frac{r_1}{\hbar\Delta_1}.
\end{equation*}
若 $N>2$，$|1\rangle$ 不直接连接另一保留端点；在完整 $Q$ 预解式的 Neumann 控制下，$r_1/(\hbar\Delta_1)=O[(\Omega_*/\Delta_*)^2]$。若 $N=2$，同一个中间态直接连接两个保留端点，必须同时保留这两个首阶源：
\[
c_1=\frac{\Omega_0c_0+\Omega_1^*c_2}{2\Delta_1}.
\]
三态开放链中这就是指定本征能量处的精确消元式。两端源可以相长或相消，却都属于同一 $O(\Omega_*/\Delta_*)$ 阶；不能把第二端源误当作下一阶。中间态总振幅受完整 $Q$ 重建算符控制，泄漏概率一般从 $O[(\Omega_*/\Delta_*)^2]$ 开始。
\end{derivation}

\begin{derivation}[反冲能量格的无量纲方程]

先保留时间依赖驱动，不立即假设矩形脉冲。二阶色散下的反冲格写成
\[
\ii\dot d_\ell(t)
=
\left[
\ell\Delta_{\rm lin}(t)
+\ell^2\omega_{\rm r}
\right]d_\ell(t)
+
\frac{\Omega(t)}2d_{\ell-1}(t)
+
\frac{\Omega^*(t)}2d_{\ell+1}(t)
+
R_\ell^{(3)}(t),
\]
其中
\[
|R_\ell^{(3)}|
\le
\epsilon_{\rm rec}^{(3)}
\,\ell^2\omega_{\rm r}|d_\ell|
\]
表示整个占据动量窗内被舍弃的三阶及更高阶色散。

取一个特征作用时间 $T$，引入无量纲时间
\[
\tau=\frac{t-t_i}{T},
\qquad
0\le\tau\le1.
\]
把 hopping 写成
\[
\Omega(t)
=
|\Omega_0|f(\tau)
\ee^{+\ii\varphi_\Omega(\tau)},
\]
并定义
\[
\varphi_{\rm lin}(\tau)
=
T\Delta_{\rm lin}(t),
\qquad
\varphi_{\rm r}
=
T\omega_{\rm r},
\qquad
\mathscr A_0
=
\frac{|\Omega_0|T}{2}.
\]
因为
\[
\partial_t
=
\frac1T\partial_\tau,
\]
原方程乘以 $T$ 后变成
\[
\ii\partial_\tau d_\ell
=
\left[
\ell\varphi_{\rm lin}(\tau)
+\ell^2\varphi_{\rm r}
\right]d_\ell
+
\mathscr A_0 f(\tau)\ee^{+\ii\varphi_\Omega(\tau)}d_{\ell-1}
+
\mathscr A_0 f(\tau)\ee^{-\ii\varphi_\Omega(\tau)}d_{\ell+1}
+
T R_\ell^{(3)}.
\]
这已经是一般单步周期驱动在二阶反冲近似下的无量纲形式。真正控制动力学的不是 $\Omega,\Delta_{\rm lin},\omega_{\rm r},T$ 四个有量纲量分别多大，而是它们组成的相位和脉冲面积。

若脉冲包络与相位在作用窗内近似常数，
\[
f(\tau)=1,
\qquad
\varphi_\Omega(\tau)=\varphi_\Omega,
\qquad
\varphi_{\rm lin}(\tau)=\varphi_{\rm lin},
\]
并且
\[
\epsilon_{\rm rec}^{(3)}
\ll1,
\]
便得到
\[
\boxed{
\ii\partial_\tau d_\ell
=
(\varphi_{\rm lin}\ell+\varphi_{\rm r}\ell^2)d_\ell
+\mathscr A_{\rm hop}\ee^{+\ii\varphi_\Omega}d_{\ell-1}
+\mathscr A_{\rm hop}\ee^{-\ii\varphi_\Omega}d_{\ell+1},
}
\]
其中
\[
\mathscr A_{\rm hop}
=
\frac{|\Omega|T}{2}.
\]
这就是正文\eqnrefs{eq:quadratic_recoil_dimensionless_lattice}。

目标边 $0\leftrightarrow1$ 的在位能差为
\[
\delta_1-\delta_0
=
\Delta_{\rm lin}+\omega_{\rm r},
\]
因此该边共振要求
\[
\Delta_{\rm lin}
=
-\omega_{\rm r},
\qquad
\varphi_{\rm lin}
=
-\varphi_{\rm r}.
\]
再看相邻两条外部边：
\[
(\delta_2-\delta_1)-(\delta_1-\delta_0)
=
2\omega_{\rm r},
\]
所以目标双态与最近外部通道之间的二阶谱间隔由
\[
2|\varphi_{\rm r}|
=
2|\omega_{\rm r}|T
\]
控制。与此同时，实际被驱动的格点宽度由 $\mathscr A_{\rm hop}$ 决定。因而二态极限不能只要求 $\varphi_{\rm r}$ 大，还必须比较
\[
\mathscr A_{\rm hop}
\ll
2|\varphi_{\rm r}|
\]
并检查更远通道；Bessel 极限则要求所有显著占据格点上的
\[
|\ell\varphi_{\rm lin}+\ell^2\varphi_{\rm r}|
\ll1.
\]
\end{derivation}
